We have already established that the DCT outperforms the DFT across all datasets. However, it is worth noting that due to the similar nature of the transforms, their performance is relatively consistent across datasets. In Fig. \ref{fig:dctvsdft}, we can see that they both perform better on the Kodak and Aerial Images datasets, and worse on Textures and Icons. This is consistent with their theory since the images in Textures and Icons are both more likely to have more rapid changes in pixel values corresponding to higher frequency waveforms.

We briefly mentioned in results how some parts of entropy coding like RLE were specifically designed for the DCT. Fig. \ref{fig:dctcrvsdcr} shows the gap between the expected compression ratio based purely on the entropy of the transformed image versus the actual compression ratio achieved. At lower quantization scales, the RLE makes little difference as there are still many nonzero values. However, as quantization is made more aggresive the returns from RLE increase massively as there are longer and longer runs of 0s.